\begin{claim}{C.70}{Proved}{Operator-compatible carriers for the original evaluation fillings}
Let $f_t$, $G$, $H$ and the $H$-invariant union $W$ of
original open $G$-orbits be as in C.64. There is one compact
path-connected family $P$ of actual global $H$-centralizing
homeomorphisms, containing the identity and every $f_t$,
preserving each original open orbit and fixing every
$G$-fixed point, with the following properties.

For a nonempty compact set $K\subset\R^2$, define its
filled hull by
\[
 \operatorname{Fill}(K)
   =\R^2\setminus U_\infty(K),
 \qquad U_\infty(K)\text{ the unbounded component of }\R^2\setminus K.
                                                               \tag{C.70.1}
\]
For $x\in W$, put $K_x=P^2(x)$. There is a unique
component $V_x$ of $\operatorname{int}\operatorname{Fill}(K_x)$
containing the entire set $P(x)$. It is simply connected.
Write $O_x$ for the original orbit containing $x$, and let
$U_x$ be the component of $V_x\cap O_x$ containing $P(x)$.
These actual domains satisfy
\[
 P(x)\subset U_x\subset O_x\cap\operatorname{Fill}(P^2(x)),
 \qquad hU_x=U_{hx}\quad(h\in H).
                                                               \tag{C.70.2}
\]
Inclusion induces an injection
\[
                  \pi_1(U_x,x)\longrightarrow\pi_1(O_x,x).
                                                               \tag{C.70.3}
\]
The original evaluation loop $t\mapsto f_t(x)$ therefore
has a based disk filling $b_x:D^2\to U_x$.
The fillings can be chosen simultaneously as a set-indexed
family satisfying the exact operator equations
\[
               b_{hx}(z)=h(b_x(z))
                    \qquad(x\in W,\ h\in H,\ z\in D^2).
                                                               \tag{C.70.4}
\]
Every choice of such fillings obeys the fixed-frontier control
\[
 \sup_{z\in D^2}|b_{x_n}(z)-y|\longrightarrow0
       \quad\text{for }x_n\to y\text{ with }y\in\operatorname{Fix}(G).
                                                               \tag{C.70.5}
\]
This is a statement about the specified individual evaluation
fillings and their simultaneous operator labels. Joint
continuity in $x$, avoidance of collisions between distinct
evaluation points, and extension to a family of homeomorphisms
are not established here.
\end{claim}
\begin{proof}
\emph{The compact family and its actual orbit neighborhoods.}
Let $\mathcal K$ be C.64's compact path-connected labelled
family, and retain its projection to actual global maps.
Put
\[
 P=\mathcal K\mathcal K^{-1}
        \ \cup\ \{f_t:0\leq t\leq1\}
        \ \cup\ \{f_t^{-1}:0\leq t\leq1\}.
                                                               \tag{C.70.6}
\]
Each of these three compact sets is path connected and
contains the identity; their union has the same properties.
Composition and inversion take place in the original global
homeomorphism group. Every displayed map commutes with each
original $h\in H$, preserves $O_x$, and fixes $\operatorname{Fix}(G)$.
The source-cover labels, when required, are the corresponding
products and inverses of C.64's retained labels and of the
original lifted loop.

By C.64.1, $x$ is in the plane interior of $P(x)$:
the orbit is an open subset of the plane. For each $p\in P$,
the image under the original homeomorphism $p$ of that
neighborhood lies in $pP(x)\subset P^2(x)$.
Consequently
\[
                   P(x)\subset\operatorname{int}P^2(x).
                                                               \tag{C.70.7}
\]
The set $P(x)$ is path connected and contains $x$.
The set $K_x=P^2(x)$ is also compact and path connected,
since it is the continuous evaluation image of $P\times P$.

\emph{Filled hulls and their simply connected interior components.}
For any compact $K$, the complement has exactly one
unbounded component: outside a sufficiently large closed
disk it is connected, and every unbounded component meets
that exterior. The filled hull is closed and bounded, hence
compact. It is contained in the Euclidean convex hull of $K$.
Indeed, a point outside that convex hull is separated from
it by a line and can be joined to infinity within the
corresponding open half-plane, without meeting $K$.

Every component $V$ of $\operatorname{int}\operatorname{Fill}(K)$
is simply connected. To prove this directly, take a Jordan
curve $J\subset V$ and its bounded Jordan disk $D_J$.
The connected unbounded set $U_\infty(K)$ is disjoint from
$J$, so it lies entirely on its unbounded side. Hence
$D_J\subset\operatorname{Fill}(K)$. The interior of
$D_J$ is open and is consequently in
$\operatorname{int}\operatorname{Fill}(K)$; it belongs
to the same component as $J$. Thus the whole disk lies in $V$.
For an arbitrary loop in $V$, polygonal approximation
inside $V$ and a finite planar regular neighborhood of
the resulting graph express it as a product of based
conjugates of its inner boundary Jordan curves. Each
bounds a disk in $V$ by what we just proved. This gives
a based null-homotopy of the original loop in $V$.

Apply this with $K=K_x$. Equation (C.70.7) places the
connected set $P(x)$ in a single component $V_x$ of
the interior of the filled hull, proving existence and
uniqueness of that component. It is simply connected
by the preceding argument. Since $P(x)\subset O_x$,
it is contained in one component $U_x$ of $V_x\cap O_x$.
This proves the inclusions in (C.70.2).

\emph{The actual inclusion on fundamental groups.}
The proof of C.65.4 works with a simply connected open
plane domain in place of its open Euclidean disk. We give
the comparison in this form. Let $O,V$ be open plane
domains with $V$ simply connected, and let $U$ be a
component of $O\cap V$. Take a based loop $\gamma$ in
$U$ which is null-homotopic in $O$.
Approximate it by a based polygonal loop within $U$,
and take a compact connected planar regular neighborhood
$M\subset U$ of its finite graph. The bounded disk
inside its outer boundary lies in $V$, by simple
connectedness and Jordan separation. In particular all
its inner complementary Jordan disks lie in $V$.

Fill precisely those inner disks wholly contained in $O$,
obtaining a compact connected planar surface $M'\subset U$.
From each remaining inner disk choose a point $z_i$ outside
$O$. The inclusion
\[
                   M'\hookrightarrow\R^2\setminus\{z_i\}
                                                               \tag{C.70.8}
\]
has the explicit strong deformation retraction of C.65.9:
choose a Schoenflies coordinate on each inner disk taking
its chosen $z_i$ to zero, apply the radial formula
$(r,\theta)\mapsto((1-\tau)r+\tau,\theta)$ there, and
apply that formula in the exterior Schoenflies coordinate
outside the outer boundary. The identity on $M'$ completes
the finite closed pasting. Every boundary and the chosen
basepoint stay fixed.
The original null-homotopy in $O$ avoids all $z_i$;
compose it with this retraction and append the polygonal
approximation homotopy. This gives a based filling in $U$.
The actual inclusion $U\hookrightarrow O$ is therefore
injective on fundamental groups. Taking $O=O_x$ and
$V=V_x$ proves (C.70.3).

By C.32, the original homeomorphism loop restricted
to $O_x$ is null-homotopic in its homeomorphism group.
Evaluation at the original point $x$ is continuous and
gives a based null-homotopy of $t\mapsto f_t(x)$ in
$O_x$. That entire loop lies in $P(x)\subset U_x$.
Equation (C.70.3) therefore gives the asserted filling
$b_x$ with image in this actual $U_x$.

\emph{All original operators and uniform control at fixed points.}
A plane homeomorphism sends the unbounded component
of the complement of a compact set to the unbounded
component of the image complement. To see unboundedness,
choose a sequence in that component escaping every
compact set. Properness of the inverse prevents its
image from lying in a compact set. It follows that
\[
          h\operatorname{Fill}(K)=\operatorname{Fill}(hK).
                                                               \tag{C.70.9}
\]
Every member of $P$ commutes with $h\in H$, so
$hP^2(x)=P^2(hx)$ and $hP(x)=P(hx)$.
Equation (C.70.9) sends the specified component $V_x$
to $V_{hx}$. Commutation with the original generated
group also gives $hO_x=O_{hx}$. Taking the indicated
components now proves $hU_x=U_{hx}$ exactly.

Choose one representative $x$ of each $H$-orbit in
$W$ and one of its based fillings in $U_x$. Define
the filling at $hx$ by $h\circ b_x$. This is independent
of the choice of $h$: an operator fixing $x$ fixes
every point of $O_x$, since each such point is $gx$
for some $g\in G$ and $h(gx)=g(hx)=gx$.
Its image therefore fixes the entire chosen filling.
The boundary identity follows from $hf_t(x)=f_t(hx)$.
This proves (C.70.4) for every original operator,
without any choice of an effective quotient representative.

Finally let $x_n\to y$ with $y\in\operatorname{Fix}(G)$.
The compact family $P^2$ fixes $y$ pointwise. Continuity
of evaluation on
$P^2\times(\{y\}\cup\{x_n:n\geq1\})$ gives
\[
             \sup_{p\in P^2}|p(x_n)-y|\longrightarrow0.
                                                               \tag{C.70.10}
\]
The Euclidean ball centered at $y$ with this radius
contains the convex hull of $P^2(x_n)$, and therefore
contains its filled hull and $U_{x_n}$. Every point of
every chosen filling lies in that same ball. This
proves (C.70.5), uniformly over the disk parameter
and independently of every filling choice.

The family just selected is equivariant as a family
indexed by points. Its proof imposes no continuity
between fillings at distinct source points. It also
imposes no configuration-space disjointness. Those
requirements, needed for an ambient homeomorphism
filling, remain part of C.12.
\end{proof}
